\documentclass[10pt,a4paper]{amsart}
\usepackage{amsmath,amssymb,amsthm,graphicx}
\usepackage{verbatim}
\newenvironment{reference}{\comment}{\endcomment}
%\newenvironment{reference}{}{}
\newenvironment{slogan}{\comment}{\endcomment}
\newenvironment{history}{\comment}{\endcomment}

% For commutative diagrams we use Xy-pic
\usepackage[all]{xy}

% We use 2cell for 2-commutative diagrams.
\xyoption{2cell}
% Theorem environments.
%
\theoremstyle{plain}
\newtheorem{theorem}[subsection]{Theorem}
\newtheorem{proposition}[subsection]{Proposition}
\newtheorem{lemma}[subsection]{Lemma}

\theoremstyle{definition}
\newtheorem{definition}[subsection]{Definition}
\newtheorem{example}[subsection]{Example}
\newtheorem{exercise}[subsection]{Exercise}
\newtheorem{situation}[subsection]{Situation}

\theoremstyle{remark}
\newtheorem{remark}[subsection]{Remark}
\newtheorem{remarks}[subsection]{Remarks}

\numberwithin{equation}{subsection}

% Macros
%
\def\lim{\mathop{\mathrm{lim}}\nolimits}
\def\colim{\mathop{\mathrm{colim}}\nolimits}
\def\Spec{\mathop{\mathrm{Spec}}}
\def\Hom{\mathop{\mathrm{Hom}}\nolimits}
\def\Ext{\mathop{\mathrm{Ext}}\nolimits}
\def\SheafHom{\mathop{\mathcal{H}\!\mathit{om}}\nolimits}
\def\SheafExt{\mathop{\mathcal{E}\!\mathit{xt}}\nolimits}
\def\Sch{\mathit{Sch}}
\def\Mor{\mathop{\mathrm{Mor}}\nolimits}
\def\Ob{\mathop{\mathrm{Ob}}\nolimits}
\def\Sh{\mathop{\mathit{Sh}}\nolimits}
\def\NL{\mathop{N\!L}\nolimits}
\def\CH{\mathop{\mathrm{CH}}\nolimits}
\def\proetale{{pro\text{-}\acute{e}tale}}
\def\etale{{\acute{e}tale}}
\def\QCoh{\mathit{QCoh}}
\def\Ker{\mathop{\mathrm{Ker}}}
\def\Im{\mathop{\mathrm{Im}}}
\def\Coker{\mathop{\mathrm{Coker}}}
\def\Coim{\mathop{\mathrm{Coim}}}

% Boxtimes
%
\DeclareMathSymbol{\boxtimes}{\mathbin}{AMSa}{"02}

%
% Macros for moduli stacks/spaces
%
\def\QCohstack{\mathcal{QC}\!\mathit{oh}}
\def\Cohstack{\mathcal{C}\!\mathit{oh}}
\def\Spacesstack{\mathcal{S}\!\mathit{paces}}
\def\Quotfunctor{\mathrm{Quot}}
\def\Hilbfunctor{\mathrm{Hilb}}
\def\Curvesstack{\mathcal{C}\!\mathit{urves}}
\def\Polarizedstack{\mathcal{P}\!\mathit{olarized}}
\def\Complexesstack{\mathcal{C}\!\mathit{omplexes}}
% \Pic is the operator that assigns to X its picard group, usage \Pic(X)
% \Picardstack_{X/B} denotes the Picard stack of X over B
% \Picardfunctor_{X/B} denotes the Picard functor of X over B
\def\Pic{\mathop{\mathrm{Pic}}\nolimits}
\def\Picardstack{\mathcal{P}\!\mathit{ic}}
\def\Picardfunctor{\mathrm{Pic}}
\def\Deformationcategory{\mathcal{D}\!\mathit{ef}}
\usepackage{cleveref}
\crefname{originalreference}{original reference}{original references}
\Crefname{originalreference}{Original reference}{Original references}
\makeatletter
\newlabel{lemma-irreducible-polynomials}{{09H0}{source}{}{sourcepointer-1}{}}
\newlabel{lemma-irreducible-polynomials@cref}{{[originalreference][0][]09H0}{[1][0][]source}{}{}{}}
\makeatother
\crefname{originalreference}{original reference}{original references}
\Crefname{originalreference}{Original reference}{Original references}
\makeatletter
\newlabel{lemma-separable-elements}{{09HC}{source}{}{sourcepointer-2}{}}
\newlabel{lemma-separable-elements@cref}{{[originalreference][0][]09HC}{[1][0][]source}{}{}{}}
\makeatother
\title{The maximal separable subextension and its purely inseparable quotient}
\author{The Stacks Project Authors}
\date{stacks-project:fields.tex:local-snapshot-20261012; GNU FDL 1.2 or later}
\begin{document}
\maketitle
\noindent\textit{Editorial note.} All delivered mathematical text is the author's own, copied from the pinned Stacks Project chapter file fields.tex; no statement or proof step was rewritten, reordered or removed.\par\medskip
\noindent\textit{Editorial note.} The counted proof is the author's complete argument: Standard separable-purely-inseparable factorization of an algebraic extension.\par\medskip
\noindent\textit{Editorial note.} The label lemma-separable-first names the counted statement, and the proof environment immediately following it in the source is delivered verbatim as its proof.\par\medskip
\noindent\textit{Editorial note.} The counted statement is carried by the author's own lemma environment in the chapter source.\par\medskip
\noindent\textit{Allowed standard prerequisites (editorial).}\begin{itemize}
\item Fields, Lemma lemma-separable-elements: the elements of an algebraic extension that are separable over the base form an intermediate field
\item Fields, Lemma lemma-irreducible-polynomials: for an algebraic element in positive characteristic there is a p-power q such that the minimal polynomial has the form P(x\textasciicircum{}q) with P separable; equivalently a suitable q-th power of the element is separable over the base
\item Purely inseparable extension: an algebraic extension in which every element has a minimal polynomial of the form x minus a q-th power of an element of the base, for some p-power q
\end{itemize}\medskip
\ifdefined\hypertarget\hypertarget{sourcepointer-1}{}\fi
\noindent\textit{Source reference.} Original reference 09H0: sources/stacks-fields/pinned-original-source.tex, source lines 1348--1348. Prerequisite pointer invoked at 1939; it names a standard fact of the same chapter that the counted argument uses, and every step of the counted proof built on it is retained verbatim in the delivered body.\par\medskip
\ifdefined\hypertarget\hypertarget{sourcepointer-2}{}\fi
\noindent\textit{Source reference.} Original reference 09HC: sources/stacks-fields/pinned-original-source.tex, source lines 1652--1652. Prerequisite pointer invoked at 1934; it names a standard fact of the same chapter that the counted argument uses, and every step of the counted proof built on it is retained verbatim in the delivered body.\par\medskip
\begin{lemma}
\label{lemma-separable-first}
\begin{slogan}
Any algebraic field extension is uniquely a separable field extension
followed by a purely inseparable one.
\end{slogan}
Let $E/F$ be an algebraic field extension. There exists a unique subextension
$E/E_{sep}/F$ such that $E_{sep}/F$ is separable and $E/E_{sep}$ is
purely inseparable.
\end{lemma}

\begin{proof}
If the characteristic is zero we set $E_{sep} = E$. Assume the characteristic
is $p > 0$. Let $E_{sep}$ be the set of elements of $E$ which are separable
over $F$. This is a subextension by Lemma \ref{lemma-separable-elements}
and of course $E_{sep}$ is separable over $F$. Given an $\alpha$ in $E$
there exists a $p$-power $q$ such that $\alpha^q$ is separable over $F$.
Namely, $q$ is that power of $p$ such that the minimal polynomial of
$\alpha$ is of the form $P(x^q)$ with $P$ separable algebraic, see
Lemma \ref{lemma-irreducible-polynomials}. Hence $E/E_{sep}$ is purely
inseparable. Uniqueness is clear.
\end{proof}

\end{document}
